\section*{Capítulo \ref{ch:g11:titration} --- A titulação}

\begin{solution}{exo:g11:titration:1}
O \omterm{def:g11:titration:titration}{titulante} é a solução de concentração conhecida com exatidão, na
bureta; a \omterm{def:g11:titration:titration}{solução titulada} contém a espécie a medir, no frasco. Na
\omterm{def:g11:titration:equivalence}{equivalência}, o \omterm{def:g11:titration:titration}{titulante} acrescentado e a espécie titulada estão nas
proporções da equação: os dois estão consumidos.
\end{solution}

\begin{solution}{exo:g11:titration:2}
$n = 0.0200 \times 0.0125 = \qty{2.50e-4}{mol}$.
\end{solution}

\begin{solution}{exo:g11:titration:3}
$V_E = \qty{14.3}{mL}$; a linha do ferro(II) começa em
\qty{14.3e-4}{mol}, isto é, \qty{1.43e-3}{mol}.
\end{solution}

\begin{solution}{exo:g11:titration:4}
O próprio \omterm{def:g11:titration:titration}{titulante} é colorido (violeta) e o seu produto é quase
incolor: a cor aparece no frasco assim que o permanganato está em
excesso.
\end{solution}

\begin{solution}{exo:g11:titration:5}
5: a equação consome 5 \omterm{def:g9:ions:ion}{íons} ferro(II) por \omterm{def:g9:ions:ion}{íon} permanganato, e na
\omterm{def:g11:titration:equivalence}{equivalência} os \omterm{def:g7:chemical-reactions:reactant}{reagentes} foram reunidos nessas proporções.
\end{solution}

\begin{solution}{exo:g11:titration:6}
$n(\ce{MnO4^-}) = 0.0200 \times 0.00840 = \qty{1.68e-4}{mol}$;
$n(\ce{Fe^{2+}}) = 5 \times \num{1.68e-4} = \qty{8.40e-4}{mol}$;
$c = \num{8.40e-4} / 0.0100 = \qty{0.0840}{mol/L}$;
$C_m = 0.0840 \times 55.8 = \qty{4.69}{g/L}$.
\end{solution}

\begin{solution}{exo:g11:titration:7}
Razão 1:1: $n = 0.0100 \times 0.0114 = \qty{1.14e-4}{mol}$ em
\qty{10.0}{mL}, logo \qty{1.14e-3}{mol} nos \qty{100.0}{mL}, isto é,
$\num{1.14e-3} \times 176.0 = \qty{0.201}{g}$: cerca de \qty{200}{mg}.
\end{solution}

\begin{solution}{exo:g11:titration:8}
Cada gota não reagiria na hora: a cor do \omterm{def:g11:titration:titration}{titulante} demoraria a sumir
antes da \omterm{def:g11:titration:equivalence}{equivalência}, e o fim seria visto cedo demais (ou seria preciso
esperar minutos depois de cada gota).
\end{solution}

\begin{solution}{exo:g11:titration:9}
Grande demais: ela passou da \omterm{def:g11:titration:equivalence}{equivalência}, acrescentando permanganato em
excesso. Perto da \omterm{def:g11:titration:equivalence}{equivalência}, ela deveria ter acrescentado o \omterm{def:g11:titration:titration}{titulante}
gota a gota, agitando, e parado no primeiro rosa claro que persiste.
\end{solution}

\begin{solution}{exo:g11:titration:10}
$n(\ce{Fe^{2+}}) = 5 \times 0.0200 \times 0.0130 = \qty{1.30e-3}{mol}$ em
\qty{10.0}{mL}: $c = \qty{0.130}{mol/L}$ para a solução diluída, e
$10 \times 0.130 = \qty{1.30}{mol/L}$ para a concentrada.
\end{solution}

\begin{solution}{exo:g11:titration:11}
$0.0630 \times 55.8 = \qty{3.52}{g}$.
\end{solution}

\begin{solution}{exo:g11:titration:12}
$0.1 / 7.2 \approx \qty{1.4}{\percent}$, contra
\qty{0.7}{\percent} para \qty{14.3}{mL}. Titular um volume maior (ou uma
amostra mais concentrada) da solução, ou usar um \omterm{def:g11:titration:titration}{titulante} mais
diluído.
\end{solution}

\begin{solution}{exo:g11:titration:13}
$n(\ce{MnO4^-}) = 0.0200 \times 0.0176 = \qty{3.52e-4}{mol}$;
$n(\ce{H2O2}) = \frac{5}{2} \times \num{3.52e-4} = \qty{8.80e-4}{mol}$ em
\qty{10.0}{mL}: \qty{0.0880}{mol/L} na solução diluída,
\qty{0.880}{mol/L} no desinfetante, isto é,
$0.880 \times 34.0 = \qty{29.9}{g/L}$ (uma solução a cerca de 3 \%).
\end{solution}

\begin{solution}{exo:g11:titration:14}
$n(\ce{MnO4^-}) = \num{1.4e-3} / 5 = \qty{2.8e-4}{mol}$, a ser fornecida
em cerca de \qty{0.015}{L}: $c \approx \qty{0.019}{mol/L}$, logo uma
solução a \qty{0.0200}{mol/L}. Dez vezes mais concentrada, $V_E$ seria
de cerca de \qty{1.4}{mL}: o erro da bureta chegaria a uns
\qty{7}{\percent}.
\end{solution}

\begin{solution}{exo:g11:titration:15}
$8 \times 0.0200 \times 0.0143 = \qty{2.29e-3}{mol}$ de \omterm{def:g9:ions:ion}{íons} hidrogênio.
\qty{10}{mL} a \qty{1.0}{mol/L} contêm \qty{1.0e-2}{mol}: bastam, com
folga de mais de quatro vezes. Sem ácido, a reação do capítulo não
poderia acontecer como está escrita (os \omterm{def:g9:ions:ion}{íons} hidrogênio são um
\omterm{def:g7:chemical-reactions:reactant}{reagente}), e o permanganato reagiria de outro modo.
\end{solution}

\begin{solution}{pb:g11:titration:1}
\textbf{1.} \ce{MnO4^-}/\ce{Mn^{2+}} e \ce{Fe^{3+}}/\ce{Fe^{2+}}. O
\omterm{def:g11:redox:oxidant}{oxidante} é o \omterm{def:g9:ions:ion}{íon} permanganato, o \omterm{def:g11:redox:oxidant}{redutor} o \omterm{def:g9:ions:ion}{íon} ferro(II).

\textbf{2.} \ce{MnO4^- + 8H+ + 5e- -> Mn^{2+} + 4H2O};
\ce{Fe^{2+} -> Fe^{3+} + e-}.

\textbf{3.} A segunda é multiplicada por 5:
\ce{MnO4^- + 8H+ + 5Fe^{2+} -> Mn^{2+} + 5Fe^{3+} + 4H2O}. Carga à
esquerda $-1 + 8 + 10 = +17$; à direita $+2 + 15 = +17$.

\textbf{4.} Os \omterm{def:g9:ions:ion}{íons} hidrogênio são um \omterm{def:g7:chemical-reactions:reactant}{reagente}: a reação precisa de uma
\omterm{def:g9:acids-bases-ph:acidic}{solução ácida}.

\textbf{5.} Ela é total, rápida e é a única reação do permanganato no
frasco; o permanganato, violeta, é consumido antes da \omterm{def:g11:titration:equivalence}{equivalência} e
colore o frasco de rosa depois dela.

\textbf{6.} Na bureta; lido na parte de baixo do menisco, com o olho na
altura, antes e depois, a \qty{0.05}{mL}.

\textbf{7.} Amarelo claro (\omterm{def:g9:ions:ion}{íons} ferro); cada gota violeta é consumida
na hora pelos \omterm{def:g9:ions:ion}{íons} ferro(II) ainda presentes.

\textbf{8.} $0.0200 \times 0.0143 = \qty{2.86e-4}{mol}$.

\textbf{9.} $n(\ce{Fe^{2+}}) = 5 \times \num{2.86e-4} =
\qty{1.43e-3}{mol}$.

\textbf{10.} $\num{1.43e-3} \times 55.8 = \qty{0.0798}{g} =
\qty{79.8}{mg}$.

\textbf{11.} $(80 - 79.8)/80 \approx \qty{0.3}{\percent}$.

\textbf{12.} $\qty{0.1}{mL}$ em \qty{14.3}{mL} é \qty{0.7}{\percent},
cerca de \qty{0.6}{mg}: o resultado, $79.8 \pm 0.6$ mg, concorda com o
rótulo.

\textbf{13.} $14.1$: $5 \times 0.0200 \times 0.0141 \times 55.8 =
\qty{78.7}{mg}$; $14.4$: \qty{80.4}{mg}.

\textbf{14.} $(79.8 + 78.7 + 80.4)/3 = \qty{79.6}{mg}$.

\textbf{15.} Os próprios comprimidos diferem um pouco no teor de ferro,
em cerca de \qty{1}{\percent}.

\textbf{16.} Não: $V_E$ seria de cerca de \qty{2.9}{mL}, e o erro da
bureta de \qty{3.5}{\percent}, cinco vezes pior.

\textbf{17.} A vitamina C também seria oxidada pelo permanganato: a
reação não seria mais única, $V_E$ ficaria grande demais, e o ferro
seria superestimado.

\textbf{18.} $\num{1.43e-3} \times 151.9 = \qty{0.217}{g}$, cerca de
\qty{217}{mg}.

\textbf{19.} $\num{1.43e-3} \times \num{6.02e23} = \num{8.6e20}$ \omterm{def:g9:ions:ion}{íons}.

\textbf{20.} Cerca de \qty{80}{mg} de ferro por comprimido: o rótulo é
honesto.
\end{solution}
